\documentclass[11pt]{article}

\usepackage[T1]{fontenc}
\usepackage[utf8]{inputenc}
\usepackage{lmodern}
\usepackage{microtype}
\usepackage[margin=1in]{geometry}
\usepackage{amsmath,amssymb,amsthm,mathtools}
\usepackage{enumitem}
\usepackage{booktabs}
\usepackage{aliascnt}
\usepackage[hidelinks,hypertexnames=false]{hyperref}
\usepackage[nameinlink,noabbrev]{cleveref}

\allowdisplaybreaks
\setlength{\parindent}{1.3em}
\setlength{\parskip}{0pt}
\emergencystretch=2em

\newtheorem{theorem}{Theorem}[section]
\newaliascnt{proposition}{theorem}
\newtheorem{proposition}[proposition]{Proposition}
\aliascntresetthe{proposition}
\newaliascnt{lemma}{theorem}
\newtheorem{lemma}[lemma]{Lemma}
\aliascntresetthe{lemma}
\newaliascnt{corollary}{theorem}
\newtheorem{corollary}[corollary]{Corollary}
\aliascntresetthe{corollary}
\theoremstyle{definition}
\newaliascnt{definition}{theorem}
\newtheorem{definition}[definition]{Definition}
\aliascntresetthe{definition}
\newaliascnt{remark}{theorem}
\newtheorem{remark}[remark]{Remark}
\aliascntresetthe{remark}

\newcommand{\Ccal}{\mathcal C}
\newcommand{\Tcal}{\mathcal T}
\newcommand{\Dcol}{\Gamma^{\mathrm d}}
\newcommand{\Fcol}{\Gamma^{\mathrm f}}
\newcommand{\set}[1]{\{#1\}}
\newcommand{\comp}{\mathbin{\setminus}}
\newcommand{\ceil}[1]{\left\lceil #1\right\rceil}
\newcommand{\target}[1]{\texttt{#1}}
\newcommand{\colset}{[k]}

\hypersetup{
  pdftitle={A complete dependency supplement for Scheide's small-Tashkinov-order theorem},
  pdfauthor={ChatGPT}
}

\title{A Complete Dependency Supplement for Scheide's\\Small-Tashkinov-Order Theorem}
\author{ChatGPT}
\date{August 2026}

\begin{document}
\maketitle

\begin{abstract}
This note supplies the proof of the result cited as \target{ThmScheide} in
Chen--Jing--Zang's Goldberg--Seymour manuscript~\cite{CJZ}:
a critical multigraph with
$\chi'(G)\ge \Delta(G)+2$ and Tashkinov order less than $11$ is elementary.
The argument follows the dependency cone of Proposition~4.5 in Scheide's
published article~\cite{S}. Numbering from Sections~2--4 of that article is
retained whenever a result is reproduced, and the labels used in the
formalization target are printed explicitly. Complete proofs are included
for Scheide's Theorem~3.3, Propositions~3.4--3.6, and every later result in
the dependency cone through Proposition~4.5. The only companion inputs are
the elementary-graph criterion and Tashkinov's elementary-tree theorem,
proved in the separately supplied file
\texttt{selected\_results\_graph\_edge\_colouring.tex}. No later or
stronger extended-Tashkinov-tree theorem is used.
\end{abstract}

\section*{Concordance with the formalization target}

The accompanying source file
\texttt{goldberg\_seymour\_trellis\_target.tex} uses the following labels.
This note deliberately reuses them, so a text or source search finds the
corresponding statement directly.

\begin{center}
\begin{tabular}{@{}lll@{}}
\toprule
Target label & Statement in this note & Source role \\
\midrule
\target{egraph2} & \Cref{egraph2} & elementary/strongly-closed criterion \\
\target{TashTree} & \Cref{TashTree} & elementarity of Tashkinov trees \\
\target{torder} & Definition on page~\pageref{torder} & Tashkinov order $t(G)$ \\
\target{ThmScheide} & \Cref{ThmScheide} & the imported small-order theorem \\
\bottomrule
\end{tabular}
\end{center}

The target calls the last statement a theorem; in Scheide's article it is
Proposition~4.5. Accordingly, \Cref{ThmScheide} below is displayed as
``Proposition~4.5'' while carrying the exact target label
\target{ThmScheide}. Bibliography keys are also chosen to agree with the
target wherever possible: in particular, the original article is
\verb|\bibitem{S}|, and the book and Tashkinov article are
\verb|\bibitem{SSTF}| and \verb|\bibitem{T}|.

The complete two-file dependency package is organized as follows. The
separately supplied book extract proves the statements represented here by
\Cref{egraph2,TashTree}; in that file they carry the searchable labels
\texttt{thm:egraph2} and \texttt{thm:tashkinov}. Starting from those two
results, the present note proves every remaining result used to derive
\target{ThmScheide}, including Scheide's Theorem~3.3 and Propositions~3.5
and~3.6. Thus the package has no additional theorem-level import in the
Scheide dependency cone.

\setcounter{section}{1}
\section{Terminology and the elementary-graph criterion}

All multigraphs are finite, loopless, and may have parallel edges. For a
multigraph $G=(V,E)$ and $X,Y\subseteq V$, let $E_G(X,Y)$ be the set of edges
with one end in $X$ and the other in $Y$. The boundary of $X$ is
\[
  \partial_G(X)=E_G(X,V\comp X).
\]
We suppress the subscript $G$ when the ambient graph is clear. The maximum
degree is $\Delta(G)$.

A proper $k$-edge-colouring of $G$ is a map $\varphi:E(G)\to\colset$ assigning
different colours to adjacent edges. The set of such colourings is
$\Ccal^k(G)$. For $v\in V$ and $X\subseteq V$, put
\[
 \varphi(v)=\set{\varphi(f):f\in E_G(v)},\qquad
 \overline\varphi(v)=\colset\comp\varphi(v),\qquad
 \overline\varphi(X)=\bigcup_{v\in X}\overline\varphi(v).
\]
The colours in $\overline\varphi(v)$ are missing at $v$. A set $X$ is
\emph{elementary} with respect to $\varphi$ if the missing-colour sets of its
distinct vertices are pairwise disjoint. It is \emph{closed} if
\[
 \varphi(\partial(X))\cap\overline\varphi(X)=\varnothing,
\]
and \emph{strongly closed} if it is closed and distinct edges of
$\partial(X)$ receive distinct colours.

For distinct colours $\alpha,\beta$, the components of the subgraph formed by
the $\alpha$- and $\beta$-coloured edges are paths or even cycles and are
called $(\alpha,\beta)$-chains. The chain containing $v$ is denoted by
$P_v(\alpha,\beta,\varphi)$. Interchanging the two colours on a chain $P$
is a Kempe change and gives the colouring $\varphi/P$. We regard an isolated
vertex as a chain of length zero. If a path $P$ has vertices $u,v$, then
$uPv$ denotes its unique subpath with ends $u,v$.

A multigraph is \emph{critical} if every proper subgraph has smaller chromatic
index. A \emph{$k$-triple} is a critical multigraph $G$ with
$\chi'(G)=k+1$ and $k\ge\Delta(G)+1$, together with an edge $e\in E(G)$ and a
colouring $\varphi\in\Ccal^k(G-e)$. The edge $e$ is regarded as uncoloured.

For completeness, let
\[
 \Gamma(G)=\max\left\{\frac{2|E_G(U,U)|}{|U|-1}:
 U\subseteq V(G),\ |U|\ge3\text{ odd}\right\}.
\]
A multigraph is \emph{elementary} if
$\chi'(G)=\ceil{\Gamma(G)}$, as in the formalization target.

\setcounter{theorem}{7}
\begin{theorem}[Elementary-graph criterion; target label \target{egraph2}]
\label{egraph2}\label{Scheide-Theorem-2.8}
Let $G=(V,E)$ be critical with $\chi'(G)=k+1$ and
$k\ge\Delta(G)+1$. The following are equivalent.
\begin{enumerate}[label=\textup{(\roman*)}]
\item $G$ is elementary.
\item For every $e\in E$ and every $\varphi\in\Ccal^k(G-e)$, the set $V$ is
      elementary with respect to $\varphi$.
\item For some $e\in E$ and some $\varphi\in\Ccal^k(G-e)$, the set $V$ is
      elementary with respect to $\varphi$.
\item For some $e\in E$, some $\varphi\in\Ccal^k(G-e)$, and some
      $X\subseteq V$ containing both ends of $e$, the set $X$ is elementary
      and strongly closed with respect to $\varphi$.
\end{enumerate}
\end{theorem}

This is the statement labelled \target{egraph2} in the target manuscript.
A complete proof is supplied in the companion file
\texttt{selected\_results\_graph\_edge\_colouring.tex}, under the label
\texttt{thm:egraph2}. Its provenance is Scheide~\cite[Theorem~2.8]{S} and
Stiebitz--Scheide--Toft--Favrholdt~\cite[Theorem~1.4]{SSTF}.

\section{Tashkinov trees and globally maximum triples}

\setcounter{theorem}{0}
\begin{definition}[Tashkinov tree; source Definition~3.1]
\label{Scheide-Definition-3.1}
Let $(G,e,\varphi)$ be a $k$-triple. A \emph{Tashkinov tree} with respect
to $e$ and $\varphi$ is a sequence
\[
 T=(y_0,e_1,y_1,\ldots,e_{n-1},y_{n-1})
\]
such that the vertices are distinct, $e_1=e$, every $e_i$ joins $y_i$ to an
earlier vertex, and, for every $2\le i\le n-1$, the colour
$\varphi(e_i)$ is missing at some $y_j$ with $j<i$. We write
$V(T)=\set{y_0,\ldots,y_{n-1}}$ and
$E(T)=\set{e_1,\ldots,e_{n-1}}$.
\end{definition}

For $1\le r\le n-1$, write
$T\mathbin{|}y_r=(y_0,e_1,y_1,\ldots,e_r,y_r)$. Every such initial segment
is again a Tashkinov tree. A colour is \emph{used on $T$} if it colours an
edge of $E(T)\comp\set e$; otherwise it is unused.

\begin{theorem}[Tashkinov; target label \target{TashTree}]
\label{TashTree}\label{Scheide-Theorem-3.2}
For every Tashkinov tree $T$ in a $k$-triple $(G,e,\varphi)$, the set $V(T)$
is elementary with respect to $\varphi$.
\end{theorem}

This is the statement labelled \target{TashTree} in the target manuscript.
A complete proof is supplied in the companion file
\texttt{selected\_results\_graph\_edge\_colouring.tex}, under the label
\texttt{thm:tashkinov}; see Tashkinov~\cite{T},
Scheide~\cite[Theorem~3.2]{S}, and
Stiebitz--Scheide--Toft--Favrholdt~\cite[Theorem~5.1]{SSTF}.

For fixed $e$ and $\varphi$, a Tashkinov tree is \emph{maximal} if no
Tashkinov tree for the same pair has a strictly larger vertex set containing
its vertex set. This fixed-pair notion is distinct from the globally maximum
notion used in the definition of $t(G)$ below.

\begin{theorem}[Maximal-tree package; source Theorem~3.3]
\label{Scheide-Theorem-3.3}\label{maximal-tree-package}
Let $(G,e,\varphi)$ be a $k$-triple. Let $T$ be maximal with respect to the
fixed pair $(e,\varphi)$, and let
$T'=(y_0,e_1,y_1,\ldots,e_n,y_n)$ be any Tashkinov tree for the same pair.
Then:
\begin{enumerate}[label=\textup{(\alph*)}]
\item $V(T)$ is elementary and closed with respect to $\varphi$;
\item $|V(T)|$ is odd;
\item $V(T')\subseteq V(T)$;
\item there is a Tashkinov tree $\widetilde T$ with
      $V(\widetilde T)=V(T)$ and
      $\widetilde T\mathbin{|}y_n=T'$;
\item if $i<j$, $\gamma\in\overline\varphi(y_i)$, and
      $\delta\in\overline\varphi(y_j)$, then $\gamma\ne\delta$ and
      $P_{y_i}(\gamma,\delta,\varphi)$ is an even path with ends
      $y_i,y_j$ and vertex set contained in $V(T)$. If $\gamma$ is unused
      on $T'\mathbin{|}y_j$, then $T'$ remains a Tashkinov tree after this
      chain is switched;
\item $G[V(T)]$ contains an odd cycle.
\end{enumerate}
\end{theorem}

\begin{proof}
Put $X=V(T)$. By \Cref{TashTree}, $X$ is elementary. It is also closed.
Indeed, if an edge $f=xz$ with $x\in X$ and $z\notin X$ had colour
$\eta\in\overline\varphi(X)$, then $f,z$ could be appended to the displayed
sequence for $T$, using a vertex of $X$ missing $\eta$ as the witness. This
would contradict maximality. This proves (a).

Every vertex of $X$ misses at least one colour. Choose
$\eta\in\overline\varphi(x)$ for some $x\in X$. Elementarity says that
$\eta$ is present at every vertex of $X\comp\set{x}$, and closedness says
that no $\eta$-edge leaves $X$. Hence the $\eta$-matching inside $X$
saturates exactly $|X|-1$ vertices. That number is even, so $|X|$ is odd.
This proves (b).

For (c), both ends of $e$ belong to $X$. If $T'$ had a vertex outside $X$,
let $y_r$ be the first such vertex in its sequence. Then $r\ge2$, all
vertices preceding $y_r$ lie in $X$, the tree edge $e_r$ joins $y_r$ to
$X$, and its colour is missing at an earlier vertex of $X$. This contradicts
closedness. Thus $V(T')\subseteq X$.

For (d), start with the sequence $T'$. While some vertex of $X$ has not yet
been used, take the first such vertex in the original sequence for $T$ and
append it with its original tree edge. Every vertex and every missing-colour
witness preceding that edge in the original sequence has already been
included. The resulting sequence is therefore a Tashkinov tree, begins with
$T'$, and eventually has vertex set $X$.

For (e), put $X'=V(T')$. The set $X'$ is elementary by
\Cref{TashTree}, so $\gamma\ne\delta$. In the
$(\gamma,\delta)$-subgraph, $y_i$ has degree one, with its incident edge
coloured $\delta$, and $y_j$ has degree one, with its incident edge coloured
$\gamma$. Every other vertex of $X$ is incident with one edge of each
colour. Since $X$ is closed and both colours are missing on $X$, no
$\gamma$- or $\delta$-edge leaves $X$. It follows that the component
through $y_i$ is a path contained in $X$ whose other end is $y_j$. Its first
edge has colour $\delta$ and its last edge has colour $\gamma$, so its
length is even.

Assume now that $\gamma$ is unused on the prefix
$T'\mathbin{|}y_j$. The colour $\delta$ is unused on that prefix as well:
if an edge $e_r$ of the prefix had colour $\delta$, the Tashkinov condition
would give an earlier vertex missing $\delta$, contrary to the elementarity
of the prefix because $\delta$ is missing at $y_j$. Thus the Kempe change
on the displayed path changes no tree edge of the prefix. At its ends it
interchanges the missing colours: afterwards $\delta$ is missing at $y_i$
and $\gamma$ at $y_j$. Every edge later than $y_j$ whose colour is
$\gamma$ or $\delta$ consequently has an earlier witness for its new or
unchanged colour at one of $y_i,y_j$; all other edge colours and witnesses
are unchanged. Hence $T'$ remains a Tashkinov tree.

Finally, let $x,y$ be the ends of $e$, in their order as the first two
vertices of $T$. Choose
$\alpha\in\overline\varphi(x)$ and
$\beta\in\overline\varphi(y)$. These colours are distinct by elementarity,
and (e) gives an even $(\alpha,\beta)$-path in $G[X]$ with ends $x,y$.
Adding the uncoloured edge $e$ produces an odd cycle in $G[X]$. This proves
(f).
\end{proof}

The statement and numbering are those of
Scheide~\cite[Theorem~3.3]{S}. The proof above uses only
\Cref{TashTree}, maximality for the fixed pair $(e,\varphi)$, and elementary
matching parity; it does not invoke a stronger extension theorem.

Let $G$ be critical with $\chi'(G)=k+1$ and $k\ge\Delta(G)+1$.
The \emph{Tashkinov order} of $G$ is
\begin{equation}
\label{torder}
 t(G)=\max |V(T)|,
\end{equation}
where the maximum ranges over all $e\in E(G)$, all
$\varphi\in\Ccal^k(G-e)$, and all corresponding Tashkinov trees. Let
$\Tcal(G)$ be the set of triples $(T,e,\varphi)$ attaining this maximum. A
triple in $\Tcal(G)$ will be called \emph{globally maximum}. Its tree is
maximal for its fixed pair, so \Cref{Scheide-Theorem-3.3} applies.

For a colour $\alpha$, put
\[
 E_\alpha(e,\varphi)=\set{f\in E(G)\comp\set e:\varphi(f)=\alpha}
\]
and
\[
 E_\alpha(T,e,\varphi)=
 E_\alpha(e,\varphi)\cap E_G(V(T),V(G)\comp V(T)).
\]
A colour $\alpha$ is \emph{defective} for $(T,e,\varphi)$ if
$|E_\alpha(T,e,\varphi)|\ge2$; the set of defective colours is
$\Dcol(T,e,\varphi)$. It is \emph{free} if it is missing on $V(T)$ and
unused on $T$; the set of free colours is $\Fcol(T,e,\varphi)$.

\begin{proposition}[Basic facts for globally maximum triples; source Proposition~3.4]
\label{Scheide-Proposition-3.4}\label{maximum-triple-facts}
Let $(T,e,\varphi)\in\Tcal(G)$. Then:
\begin{enumerate}[label=\textup{(\alph*)}]
\item $t(G)=|V(T)|\ge3$ is odd;
\item $V(T)$ is elementary and closed with respect to $\varphi$;
\item if $\Dcol(T,e,\varphi)=\varnothing$, then $V(T)$ is strongly closed;
\item if $\alpha\in\overline\varphi(V(T))$, then
      $E_\alpha(T,e,\varphi)=\varnothing$;
\item if $\alpha\in\Dcol(T,e,\varphi)$, then
      $|E_\alpha(T,e,\varphi)|\ge3$ is odd;
\item for $v\in V(T)$,
\[
 |\overline\varphi(v)|=
 \begin{cases}
 k-d_G(v)+1\ge2,&\text{if $v$ is an end of $e$},\\
 k-d_G(v)\ge1,&\text{otherwise};
 \end{cases}
\]
\item $|\Fcol(T,e,\varphi)|\ge4$;
\item every colour in
      $\Dcol(T,e,\varphi)\cup\Fcol(T,e,\varphi)$ is unused on $T$.
\end{enumerate}
\end{proposition}

\begin{proof}
Parts (a) and (b) follow from
\Cref{Scheide-Theorem-3.3}(a),(b): a maximal tree contains the initial
2-vertex tree and has odd order, hence has at least three vertices. Part (c)
is the definition of strongly closed once there is no defective boundary
colour, and (d) is closedness.

For (e), a defective colour $\alpha$ occurs on a boundary edge, so closedness
implies that it is present at every vertex of $V(T)$. The
$\alpha$-coloured edges form a matching. Counting incidences at the odd set
$V(T)$ gives
\[
 |V(T)|=2|E_\alpha(e,\varphi)\cap E_G(V(T),V(T))|
       +|E_\alpha(T,e,\varphi)|.
\]
Thus the boundary number is odd; if it is at least two, it is at least three.
Part (f) is the degree count in $G-e$.

By elementarity and (f), the union of the missing-colour sets has size at
least $|V(T)|+2$. At most $|V(T)|-2$ colours occur on the coloured tree
edges. Every used tree-edge colour is missing at an earlier tree vertex, so
at least four missing colours are unused. This proves (g). Finally, free
colours are unused by definition. A defective colour is present at every
tree vertex, whereas every used tree-edge colour is missing at an earlier
one; hence a defective colour is also unused. This proves (h).
\end{proof}

A vertex $v\in V(G)\comp V(T)$ is \emph{absorbing} with respect to
$(T,e,\varphi)$ if, for every $\delta\in\overline\varphi(v)$ and every
$\gamma\in\Fcol(T,e,\varphi)\comp\set\delta$, the chain
$P_v(\gamma,\delta,\varphi)$ contains the unique vertex
$u\in V(T)$ for which $\gamma\in\overline\varphi(u)$. In that event the
chain is a path with ends $u,v$. Let $A(T,e,\varphi)$ be the set of
absorbing vertices.

\begin{proposition}[Absorbing vertices; source Proposition~3.5]
\label{Scheide-Proposition-3.5}\label{absorbing-elementary}
For every $(T,e,\varphi)\in\Tcal(G)$, the set
$V(T)\cup A(T,e,\varphi)$ is elementary with respect to $\varphi$.
\end{proposition}

\begin{proof}
Put $X=V(T)$ and $A=A(T,e,\varphi)$. The set $X$ is elementary by
\Cref{maximum-triple-facts}(b). We first show that a colour missing at a
vertex of $A$ is missing at no vertex of $X$. Suppose that
$v\in A$, $\delta\in\overline\varphi(v)$, and
$\delta\in\overline\varphi(x)$ for some $x\in X$. Choose
$\gamma\in\Fcol(T,e,\varphi)\comp\set{\delta}$, possible by
\Cref{maximum-triple-facts}(g), and let $u\in X$ be the unique vertex
missing $\gamma$. Both $\gamma$ and $\delta$ are missing on the closed set $X$, so no
edge of either colour leaves $X$. Consequently the
$(\gamma,\delta)$-component containing the outside vertex $v$ cannot meet
$X$, contrary to the absorbing property, which requires it to contain $u$.

It remains to compare two distinct vertices $v,w\in A$. Suppose that a
colour $\delta$ is missing at both. Choose
$\gamma\in\Fcol(T,e,\varphi)\comp\set{\delta}$ and let $u\in X$ be the
vertex missing $\gamma$. By the first paragraph, $\delta$ is present at
$u$. The absorbing property says that the
$(\gamma,\delta)$-component containing $u$ is a path with ends $u,v$; the
same component is also a path with ends $u,w$. Since $v\ne w$, this is
impossible. Thus the missing-colour sets of all distinct vertices in
$X\cup A$ are pairwise disjoint.
\end{proof}

This is the absorbing-vertex theorem quoted in
Scheide~\cite[Proposition~3.5]{S}; its earlier source is
Favrholdt--Stiebitz--Toft~\cite{FST}. The proof is included here so that no
result beyond the two companion inputs remains implicit.

Fix an end $w$ of a path $P$. Write $x\preceq_{w,P}y$ when $x$ lies on the
subpath $wPy$.

\begin{proposition}[Defective-chain structure; source Proposition~3.6]
\label{Scheide-Proposition-3.6}\label{defective-chain}
Let $(T,e,\varphi)\in\Tcal(G)$, let
$\alpha\in\Dcol(T,e,\varphi)$, let
$\gamma\in\Fcol(T,e,\varphi)$ be missing at $u\in V(T)$, and put
$P=P_u(\alpha,\gamma,\varphi)$. Then:
\begin{enumerate}[label=\textup{(\alph*)}]
\item $P$ is a path from $u$ to a vertex outside $V(T)$;
\item
$E_\alpha(T,e,\varphi)=
 E(P)\cap E_G(V(T),V(G)\comp V(T))$;
\item in the order $\preceq_{u,P}$ there is a first vertex $v_1$ outside
      $V(T)$ and a last vertex $v_2$ in $V(T)$ after that first departure;
\item $\overline\varphi(v_2)\cap\Fcol(T,e,\varphi)=\varnothing$;
\item $V(T)\cup\set{v_1}$ is elementary with respect to $\varphi$.
\end{enumerate}
\end{proposition}

\begin{proof}
Put $X=V(T)$. By
\Cref{maximum-triple-facts}, $X$ is elementary and closed, every defective
or free colour is unused on $T$, $\alpha$ is present at every vertex of
$X$, no $\gamma$-edge leaves $X$, and at least three
$\alpha$-edges leave $X$.

At $u$ the $(\alpha,\gamma)$-subgraph has degree one, whereas every other
vertex of $X$ has degree two. Thus the component $P$ through $u$ is a path,
and its other end lies outside $X$. This proves (a).

We prove (b) by contradiction. Let $f$ be an $\alpha$-coloured boundary
edge not on $P$, and let $Q$ be its $(\alpha,\gamma)$-component. The
component $Q$ does not contain $u$. It contains no tree edge, since both
colours are unused on $T$, and it has no end in $X$: every vertex of
$X\comp\set{u}$ is incident with both colours. Switch $Q$. The displayed
sequence for $T$ is still Tashkinov, and $\gamma$ is still missing at $u$.
The edge $f$ now has colour $\gamma$, so it and its outside end can be
appended to $T$, contradicting the definition of $t(G)$. Hence every
$\alpha$-coloured boundary edge lies on $P$.

The path starts in $X$ and ends outside it. Its boundary crossings all have
colour $\alpha$, because no $\gamma$-edge leaves $X$. By (b) there are at
least three such crossings. Consequently the path returns to $X$ after its
first departure, and the vertices $v_1,v_2$ in (c) exist.

For (d), let $z$ be the outside end of $P$ and put $R=v_2Pz$. Then
$V(R)\cap X=\set{v_2}$ and $R$ contains exactly one boundary edge of $X$.
Suppose that a free colour $\beta$ is missing at $v_2$. Since $v_2$ occurs
after the first departure, it is distinct from $u$, so elementarity gives
$\beta\ne\gamma$. Interchange the names $\beta$ and $\gamma$ on all edges
with both ends outside $X$. This is a proper colouring and leaves every
edge incident with $X$ unchanged, because both colours are missing on the
closed set $X$. Under the resulting colouring, $R$ is the full
$(\alpha,\beta)$-chain from $v_2$. Switch it. No tree edge changes, and the
only missing-colour change on $X$ replaces $\beta$ by $\alpha$ at $v_2$;
since both colours are unused on $T$, the tree remains Tashkinov. Now
$\alpha$ is missing at $v_2$, while at least two of the original
$\alpha$-coloured boundary edges lie outside $R$ and retain their colour.
Appending either one contradicts the maximality of $t(G)$. Thus (d) holds.

For (e), suppose that a colour $\delta$ is missing both at $v_1$ and at
some $x\in X$. The vertex $v_1$ is not an end of $P$: by (b) the path has
at least two further boundary crossings after its first departure. Hence
both $\alpha$ and $\gamma$ are present at $v_1$, and
$\delta\notin\set{\alpha,\gamma}$. Let
$C=P_{v_1}(\gamma,\delta,\varphi)$. Both colours are missing on the closed
set $X$, so no edge of either colour leaves $X$. Therefore $C$, which
contains the outside vertex $v_1$, is disjoint from $X$.
Switch $C$. The tree remains Tashkinov, $\gamma$ remains missing at $u$,
and $\gamma$ is now missing at $v_1$. Moreover, the edge of $P$ following
$v_1$ was the unique $\gamma$-edge at $v_1$ and has changed to
$\delta$. Thus the initial segment $uPv_1$ is now the full
$(\alpha,\gamma)$-chain from $u$ to $v_1$.

Switch this initial segment. The colours $\alpha,\gamma$ are unused on
$T$, so the tree remains Tashkinov, and $\alpha$ becomes missing at $u$.
The segment contains only the first $\alpha$-coloured boundary edge of
$P$; at least two further such edges are unchanged. Appending one of them
again contradicts $t(G)$. Therefore no colour missing at $v_1$ is missing
on $X$, and $X\cup\set{v_1}$ is elementary.
\end{proof}

This is the chain theorem quoted in
Scheide~\cite[Proposition~3.6]{S}, originally due to
Favrholdt--Stiebitz--Toft~\cite{FST}. Its proof above uses only the
maximal-tree package already proved in this note.

Let $u\in V(T)$, let $\gamma\in\overline\varphi(u)$, and let
$\delta\in\Dcol(T,e,\varphi)$. The chain
$P_u(\gamma,\delta,\varphi)$ is a path from $u$ to a vertex outside $V(T)$:
$\delta$ is present at every tree vertex, and elementarity leaves no second
tree vertex missing $\gamma$. Its last vertex in $V(T)$, in the order
starting at $u$, is an \emph{exit vertex}. Let $F(T,e,\varphi)$ denote the
set of all exit vertices.

\setcounter{theorem}{7}
\begin{lemma}[Exit vertices miss no free colour; source Lemma~3.8]
\label{Scheide-Lemma-3.8}\label{exit-no-free}
For every $(T,e,\varphi)\in\Tcal(G)$,
\[
 \overline\varphi(F(T,e,\varphi))
 \cap\Fcol(T,e,\varphi)=\varnothing.
\]
\end{lemma}

\begin{proof}
Let $v$ be an exit vertex. Thus, for some $u\in V(T)$,
$\gamma\in\overline\varphi(u)$, and
$\delta\in\Dcol(T,e,\varphi)$, the vertex $v$ is the last tree vertex on
$P=P_u(\gamma,\delta,\varphi)$; the other end $z$ of $P$ lies outside
$V(T)$. Suppose that a free colour $\alpha$ is missing at $v$.

If $\alpha=\gamma$, elementarity forces $u=v$. Then
\Cref{defective-chain}(d), applied directly to the free colour $\alpha$ and
the defective colour $\delta$, says that the last tree vertex on this chain
misses no free colour, a contradiction.

Assume $\alpha\ne\gamma$. Since both colours are missing on the closed set
$V(T)$, neither appears on its boundary. Interchange $\alpha$ and $\gamma$
on every edge whose two ends lie outside $V(T)$. The resulting colouring
$\varphi'$ agrees with $\varphi$ on every edge incident with $V(T)$; hence
$(T,e,\varphi')\in\Tcal(G)$ and its free and defective colour sets are
unchanged. The subpath $vPz$ is now the
$(\alpha,\delta)$-chain from $v$, and it crosses the boundary of $V(T)$
exactly once. This contradicts \Cref{defective-chain}(b), which requires
that chain to contain every $\delta$-coloured boundary edge, of which there
are at least three. Therefore no free colour is missing at $v$.
\end{proof}

\subsection*{Normal and balanced globally maximum trees}

Let $(T,e,\varphi)\in\Tcal(G)$, where
\[
 T=(y_0,e_1,y_1,\ldots,e_{n-1},y_{n-1}),\qquad n=t(G).
\]
The tree is \emph{normal} if there are
$\alpha\in\overline\varphi(y_0)$,
$\beta\in\overline\varphi(y_1)$, an integer $2\le p\le n$, and an edge
$q\in E_G(y_{p-1},y_0)$ such that
\[
 P(y_1,e_2,y_2,\ldots,e_{p-1},y_{p-1},q,y_0)
\]
is an $(\alpha,\beta)$-chain. The initial tree
$T\mathbin{|}y_{p-1}$ is its $(\alpha,\beta)$-\emph{trunk}, and its height
is $h(T)=p$. Let $\Tcal^N(G)$ be the set of normal globally maximum triples,
and let $h(G)$ be the maximum height of a member of $\Tcal^N(G)$.

\begin{remark}[Indexing convention]
The upper bound is $p\le n$, so the trunk is allowed to contain all
$n=t(G)$ vertices. This is the convention required in source
Proposition~4.2, whose first case is $h(G)=t(G)$. The bound $p\le n-1$
printed in the 2008 preprint is therefore read as an indexing slip.
\end{remark}

\setcounter{theorem}{10}
\begin{lemma}[Normalisation; source Lemma~3.11]
\label{Scheide-Lemma-3.11}\label{normalisation}
Let $(T,e,\varphi)\in\Tcal(G)$ and let $e=xy$. There are
$\alpha\in\overline\varphi(x)$ and
$\beta\in\overline\varphi(y)$ and a normal Tashkinov tree $T'$ with
$V(T')=V(T)$ such that
\[
 h(T')=|V(P_x(\alpha,\beta,\varphi))|.
\]
\end{lemma}

\begin{proof}
Choose $\alpha\in\overline\varphi(x)$ and
$\beta\in\overline\varphi(y)$, possible because the degree of each end of
the uncoloured edge in $G-e$ is at most $\Delta(G)-1<k$. By
\Cref{Scheide-Theorem-3.3}(e), the colours are distinct and
$P=P_x(\alpha,\beta,\varphi)$ is an even path with ends $x,y$ and all
vertices in $V(T)$. Write
\[
 P=P(v_1,f_2,v_2,\ldots,f_p,v_p),\qquad v_1=y,\quad v_p=x.
\]
Then
\[
 (x,e,y,f_2,v_2,\ldots,f_{p-1},v_{p-1})
\]
is a Tashkinov tree: every added edge has colour $\alpha$ or $\beta$, and
each of those colours is missing at one of the first two vertices. By
\Cref{Scheide-Theorem-3.3}(d), this initial tree extends to a Tashkinov tree
on the full vertex set $V(T)$. It is normal and has height $|V(P)|$.
\end{proof}

A normal triple $(T,e,\varphi)$ is \emph{balanced} if
$h(T)=h(G)=p$ and
\[
 \varphi(e_{2j-1})=\varphi(e_{2j})
 \qquad\text{whenever }p<2j<t(G).
\]
Let $\Tcal^B(G)$ be the set of balanced triples.

\begin{lemma}[Balanced triples; source Lemma~3.12]
\label{Scheide-Lemma-3.12}\label{balanced-triples}
The following hold.
\begin{enumerate}[label=\textup{(\alph*)}]
\item $h(G)\ge3$ is odd.
\item If $(T,e,\varphi)\in\Tcal^N(G)$ and $h(T)=h(G)$, then there is a
      Tashkinov tree $T'$ on the same vertex set for which
      $(T',e,\varphi)\in\Tcal^B(G)$; every colour used on $T'$ is used on
      $T$.
\item $\Tcal^B(G)\ne\varnothing$.
\end{enumerate}
\end{lemma}

\begin{proof}
Let $p=h(G)$ and let $T\mathbin{|}y_{p-1}$ be an
$(\alpha,\beta)$-trunk. Its closing $(\alpha,\beta)$-path has distinct ends
$y_0,y_1$, and exactly one of $\alpha,\beta$ is missing at each end. It has
positive even length, so the number $p$ of its vertices is odd and at least
three. This proves (a).

For (b), fix a normal tree of height $p$. Among all Tashkinov trees $S$ for
the same pair $(e,\varphi)$ that have the prescribed trunk, use only colours
used on the original tree, and have their post-trunk edges in consecutive
equal-coloured pairs, choose one with a maximum odd number $i$ of vertices.
Initially $i\ge p$. Suppose $i<t(G)$. Let $y_r$ be the first vertex of the
original tree not in $S$, and append the original edge $e_r$ and $y_r$.
Its colour $\gamma$ is missing at an earlier vertex, so the enlarged
Tashkinov tree has even order and exactly one of its vertices misses
$\gamma$, by \Cref{TashTree}. Counting incidences of the
$\gamma$-matching shows that an odd, hence nonzero, number of
$\gamma$-edges leave this even vertex set. Append one such edge and its
outside end. The two newly appended edges have the same colour, so all
requirements are preserved and the odd order has increased by two, a
contradiction. Hence $i=t(G)$. Applying
\Cref{Scheide-Theorem-3.3}(c),(d) to the original globally maximum tree, the
resulting tree has the same vertex set. This proves (b), and (c) follows
from \Cref{normalisation}.
\end{proof}

We shall repeatedly rotate a balanced trunk. Let
$(T,e,\varphi)\in\Tcal^B(G)$, let $p=h(G)$, set $f_1=e$ and
$f_i=e_i$ for $2\le i\le p-1$, and let
$f_p\in E_G(y_{p-1},y_0)$ be the closing edge. The cycle
$f_1,\ldots,f_p$ has $f_1$ uncoloured and the remaining edges alternately
coloured $\alpha,\beta$. For $1\le j\le p-1$, move the uncoloured status
successively around this cycle from $f_1$ to $f_{j+1}$, shifting the two
colours on the intervening edges. This gives a colouring
$\varphi_j\in\Ccal^k(G-f_{j+1})$ agreeing with $\varphi$ off the cycle and a
new trunk
\[
 (y_j,f_{j+1},y_{j+1},\ldots,f_p,y_0,
   f_1,y_1,\ldots,f_{j-1},y_{j-1}).
\]
Appending the unchanged tail produces another balanced globally maximum
triple, denoted by
\[
 (T,e,\varphi)(y_0\to y_j).
\]
Thus any trunk vertex may be made the first vertex without changing the
underlying vertex set or the tail blocks.


\section{The small-Tashkinov-order theorem}

The next lemma is the structural core of Scheide's argument.  We retain the
source hypotheses and the source's numbered intermediate conclusions.  In
particular, the proof below is not an appeal to the original Lemma~4.1; it is
a reconstruction of that proof with every recolouring and parity step
specified.

\begin{lemma}[Structural lemma; source Lemma~4.1]
\label{Scheide-Lemma-4.1}\label{structural-lemma}
Let $G$ be critical with $\chi'(G)=k+1$ and $k\ge\Delta(G)+1$, and let
$(T,e,\varphi)\in\Tcal(G)$. Suppose that
\[
\begin{split}
T={}&(y_0,e_1,y_1,\ldots,e_{p-1},y_{p-1},
 f^1_{\gamma_1},u^1_{\gamma_1},f^2_{\gamma_1},u^2_{\gamma_1},\ldots,\\
&\hspace{42mm}
 f^1_{\gamma_s},u^1_{\gamma_s},f^2_{\gamma_s},u^2_{\gamma_s}),
\end{split}
\]
where $\Gamma=\set{\gamma_1,\ldots,\gamma_s}$ and
$Y=\set{y_0,\ldots,y_{p-1}}$, and suppose that:
\begin{enumerate}[label=\textup{(S\arabic*)},leftmargin=3.2em]
\item $f^j_\gamma\in E_G(Y,u^j_\gamma)$ for every
      $\gamma\in\Gamma$ and $j\in\set{1,2}$;
\item
      $\varphi(f^1_\gamma)=\varphi(f^2_\gamma)=\gamma$ and
      $\gamma\in\overline\varphi(Y)$ for every $\gamma\in\Gamma$;
\item $\varphi(e_j)\notin\Gamma$ for $2\le j\le p-1$;
\item every vertex of $F(T,e,\varphi)$ misses a colour of $\Gamma$.
\end{enumerate}
Then $G$ is elementary.
\end{lemma}

\begin{proof}
By \Cref{maximum-triple-facts}(b), $V(T)$ is elementary and closed. If
$\Dcol(T,e,\varphi)=\varnothing$, it is strongly closed by
\Cref{maximum-triple-facts}(c), and \Cref{egraph2}(iv) proves the result.
We therefore assume throughout that
$\Dcol(T,e,\varphi)\ne\varnothing$.

We first reduce the data without changing the vertex set or the colouring.
Each colour of $\Gamma$ is missing at a unique vertex of $V(T)$. By (S4),
choose one such colour for each exit vertex and call the resulting set
$\Gamma'$. Distinct exit vertices yield distinct colours, by elementarity,
so $|\Gamma'|=|F(T,e,\varphi)|$. Reorder the two-vertex blocks so that the
blocks with colours in $\Gamma\comp\Gamma'$ precede those with colours in
$\Gamma'$, and enlarge $Y$ by the vertices of the former blocks. This
reordering is legitimate because both edges of every block join its two new
vertices to the original set $Y$ and their common colour is missing on that
set. The four hypotheses remain true with $\Gamma'$ and the enlarged $Y$:
for (S3), the newly absorbed block colours lie outside $\Gamma'$. We may
therefore replace the data by these reduced data and assume
\begin{equation}
 |F(T,e,\varphi)|=s,
 \qquad
 \text{each exit vertex misses exactly one colour of $\Gamma$.}
 \tag{1}\label{L41-1}
\end{equation}
The map sending an exit vertex to its missing colour in $\Gamma$ is
injective and both sets have size $s$; hence it is a bijection. In
particular, every colour of $\Gamma$ is missing at an exit vertex.

Every colour used on $T$ is missing on $Y$. For the initial segment this is
the Tashkinov condition, and for the remaining edges it is (S2). Every
vertex of $V(T)\comp Y$ misses at least one colour. Such a colour cannot be
used on $T$, since a used colour is already missing at a vertex of $Y$ and
$V(T)$ is elementary. It is therefore free. By \Cref{exit-no-free},
\begin{equation}
 F(T,e,\varphi)\subseteq Y.
 \tag{2}\label{L41-2}
\end{equation}

For $\gamma\in\Gamma$, put
\[
 U_\gamma=\set{u^1_\gamma,u^2_\gamma},
 \qquad
 Z_\gamma=V(T)\comp U_\gamma,
\]
and let $T-\gamma$ be the Tashkinov tree obtained by deleting the two
$\gamma$-edges and their two new vertices. We claim that
\begin{equation}
 E_G(u^1_\gamma,u^2_\gamma)\cap E_\alpha(e,\varphi)
 \ne\varnothing
 \quad
 \text{for every }\alpha\in\overline\varphi(Z_\gamma)
                 \comp\set{\gamma}.
 \tag{3}\label{L41-3}
\end{equation}
Indeed, $\alpha$ is missing on the closed set $V(T)$, so the unique
$\alpha$-edge at $u^1_\gamma$ has its other end in $V(T)$. Suppose that
this other end lies in $Z_\gamma$. The unique $\alpha$-edge at
$u^2_\gamma$ cannot end at $u^1_\gamma$, which is already incident with
an $\alpha$-edge, and therefore also ends in $Z_\gamma$. Appending these
two $\alpha$-edges to $T-\gamma$ gives another globally maximum Tashkinov
tree with vertex set $V(T)$. In that tree $\gamma$ is missing and unused,
so it is free. Its defective colours and exit vertices are unchanged,
because these depend only on the colouring and the vertex set. By the
bijection following \eqref{L41-1}, an exit vertex misses $\gamma$, contrary
to \Cref{exit-no-free}. Thus the $\alpha$-edge at $u^1_\gamma$ must join
the two vertices of $U_\gamma$, proving \eqref{L41-3}.

Both ends of the uncoloured edge $e$ lie in $Y$. By
\Cref{maximum-triple-facts}(f), $|\overline\varphi(y_0)|\ge2$. Moreover,
$y_0$ misses at most one colour of $\Gamma$: every colour in $\Gamma$ is
missing at an exit vertex, and \eqref{L41-1} says that an exit vertex
misses only one such colour. Fix
\[
 \alpha_0\in\overline\varphi(y_0)\comp\Gamma.
\]
For $\delta\in\Dcol(T,e,\varphi)$, let
\[
 P_\delta=P_{y_0}(\alpha_0,\delta,\varphi).
\]
Since $\alpha_0$ is missing on $V(T)$ and the defective colour $\delta$ is
present at every vertex of $V(T)$, this is a path from $y_0$ to a vertex
$z_\delta\notin V(T)$. Let $v^0_\delta$ be its last vertex in $V(T)$,
and let $f^0_\delta$ be the following edge, with outside end
$u^0_\delta$. The edge $f^0_\delta$ has colour $\delta$, because no
$\alpha_0$-edge leaves the closed set $V(T)$. The vertex $v^0_\delta$ is
an exit vertex, so by \eqref{L41-1} there is a unique colour
$\gamma(\delta)\in\Gamma$ missing at $v^0_\delta$.

We next prove
\begin{equation}
 E_G(U_{\gamma(\delta)},Z_{\gamma(\delta)})
       \cap E_\delta(e,\varphi)=\varnothing
 \quad\text{for every }\delta\in\Dcol(T,e,\varphi).
 \tag{4}\label{L41-4}
\end{equation}
Suppose instead that a $\delta$-edge $g_1$ joins
$u^1_{\gamma(\delta)}$ to $Z_{\gamma(\delta)}$ (renumber the two vertices
of the pair if necessary). By \eqref{L41-3}, an $\alpha_0$-edge $g_2$
joins the two vertices of $U_{\gamma(\delta)}$. Since $g_1$ occupies the
$\delta$-incidence at one of these vertices, at most one
$\delta$-coloured boundary edge of $V(T)$ is incident with the pair. A
defective colour has at least three boundary edges by
\Cref{maximum-triple-facts}(e), so there is a $\delta$-edge
$g_3\ne f^0_\delta$ from $Z_{\gamma(\delta)}$ to a vertex outside
$V(T)$.

Interchange $\alpha_0$ and $\gamma(\delta)$ on every edge with both ends
outside $V(T)$. This is a proper colouring because neither colour occurs
on $\partial(V(T))$. Under this colouring the subpath
$v^0_\delta P_\delta z_\delta$ is a
$(\gamma(\delta),\delta)$-chain. Switch that subpath. The tree
$T-\gamma(\delta)$ remains Tashkinov: no one of its tree edges changes,
$\delta$ is unused on $T$, no $\gamma(\delta)$-edge remains in the tree,
and the only changed missing status on its vertex set replaces
$\gamma(\delta)$ by $\delta$ at $v^0_\delta$. Thus $\delta$ is now missing
at $v^0_\delta$. The three edges $g_1,g_2,g_3$ retain, respectively, the
colours $\delta,\alpha_0,\delta$. Appending them successively adds
$u^1_{\gamma(\delta)}$, $u^2_{\gamma(\delta)}$, and a vertex outside
$V(T)$ after only the two vertices of $U_{\gamma(\delta)}$ were deleted.
This produces a Tashkinov tree with more than $t(G)$ vertices, a
contradiction. Hence \eqref{L41-4} holds.

Put
\[
 P'_\delta=P_{v^0_\delta}(\gamma(\delta),\delta,\varphi).
\]
It is a path whose other end $z'_\delta$ lies outside $V(T)$. We claim
\begin{equation}
 V(P'_\delta)\cap V(T)=\set{v^0_\delta}.
 \tag{5}\label{L41-5}
\end{equation}
Suppose not. Following $P'_\delta$ from $v^0_\delta$, let
$v_1\ne v^0_\delta$ be its last vertex in $V(T)$. Then $v_1$ is an exit
vertex and lies in $Y$ by \eqref{L41-2}. Let $f_1$ be the
$\delta$-edge by which the path leaves $v_1$, and let $u_1$ be its outside
end. Write $\gamma=\gamma(\delta)$, and let
$\gamma_j\in\Gamma$ be the unique colour missing at $v_1$. Since
$v_1\ne v^0_\delta$, elementarity gives $\gamma_j\ne\gamma$.

First interchange $\alpha_0$ and $\gamma$ outside $V(T)$ and then switch
$v^0_\delta P_\delta z_\delta$. As in the proof of \eqref{L41-4}, this
gives a colouring $\varphi_2$ under which $T-\gamma$ is Tashkinov,
$\delta$ is missing at $v^0_\delta$, and $f_1$ still has colour
$\delta$. Append $f_1,u_1$. Independently, interchange $\gamma$ and
$\gamma_j$ outside $V(T)$ and then switch
$v_1P'_\delta z'_\delta$. This gives a colouring $\varphi'_2$ under
which $T-\gamma_j$ is Tashkinov, $\delta$ is missing at $v_1$, and
$f^0_\delta$ is unchanged. Append $f^0_\delta,u^0_\delta$.

Let
\[
 Z=V(T)\comp(U_\gamma\cup U_{\gamma_j}).
\]
The set $Z$ has odd order, and $\delta$ is present at every vertex of $Z$.
Counting incidences of the $\delta$-matching shows that an odd number of
$\delta$-edges leave $Z$. Two such edges are $f^0_\delta$ and $f_1$;
therefore there is a third one, say $f_2$, with other end $u_2$. This edge
is unchanged by both recolourings: it has an end in $Z$, and, being
distinct from $f^0_\delta,f_1$, it lies on neither switched tail.

The vertex $u_2$ is not in $U_{\gamma_j}$. Otherwise, after appending
$f_2,u_2$ to the second partial tree, \eqref{L41-3} supplies an internal
$\alpha_0$-edge of $U_{\gamma_j}$, and appending that edge and the other
vertex of the pair gives more than $t(G)$ vertices. Nor is $u_2$ in
$U_\gamma$, by \eqref{L41-4}. Hence $u_2\notin V(T)$. Append
$f_2,u_2$ to both partial trees. Each resulting tree has exactly $t(G)$
vertices and is therefore globally maximum.

By \Cref{maximum-triple-facts}(f) and elementarity, at least four distinct
colours are missing on $\set{y_0,y_1}$. Choose one, $\beta$, outside
$\set{\alpha_0,\gamma,\gamma_j}$. It is also distinct from $\delta$,
which is present throughout $V(T)$. None of the preceding recolourings
changes a $\beta$-edge or the fact that $\beta$ is missing at an original
tree vertex. In the first new globally maximum tree, $\beta$ is therefore
present at $u_2$ and, by closedness, its edge at $u_2$ has its other end in
that tree. The original set $V(T)$ is closed and misses $\beta$, so this
other end cannot lie in $V(T)$; it must be $u_1$. In the second new
globally maximum tree the vertex $u_1$ is absent, while the same
$\beta$-edge leaves the tree and $\beta$ is missing on it. Appending that
edge and $u_1$ yields more than $t(G)$ vertices, a contradiction. This
proves \eqref{L41-5}.

For each defective colour $\delta$, define
\[
 \varphi_\delta=\varphi/P'_\delta.
\]
By \eqref{L41-5}, this switch has the following precisely delimited effect:
\begin{equation}
\begin{gathered}
 \varphi_\delta(f^0_\delta)=\gamma(\delta),\qquad
 \overline{\varphi_\delta}(v^0_\delta)=
 \bigl(\overline\varphi(v^0_\delta)\comp\set{\gamma(\delta)}\bigr)
       \cup\set{\delta},\\
 \varphi_\delta(f)=\varphi(f)
 \quad\text{for every other edge $f$ incident with $V(T)$},\\
 \overline{\varphi_\delta}(v)=\overline\varphi(v)
 \quad\text{for }v\in V(T)\comp\set{v^0_\delta},
 \quad\text{and for every }v\notin V(P'_\delta).
\end{gathered}
 \tag{6}\label{L41-6}
\end{equation}

We next claim
\begin{equation}
 \left|E_G(Z_{\gamma(\delta)},V(G)\comp V(T))
       \cap E_\delta(e,\varphi)\right|=3
 \quad\text{for every }\delta\in\Dcol(T,e,\varphi).
 \tag{7}\label{L41-7}
\end{equation}
The displayed number is odd. Indeed, $Z_{\gamma(\delta)}$ has odd order,
$\delta$ is present at every one of its vertices, and \eqref{L41-4}
excludes $\delta$-edges from it to $U_{\gamma(\delta)}$. If the number were
greater than three, then, besides $f^0_\delta$, three of these edges would
retain colour $\delta$ under $\varphi_\delta$ by \eqref{L41-6}. Appending
their three distinct outside ends to $T-\gamma(\delta)$ would give
$t(G)+1$ vertices.

Suppose instead that the number is less than three. It is then one, and its
only edge is $f^0_\delta$. By \eqref{L41-3}, the two vertices of
$U_{\gamma(\delta)}$ are joined by an $\alpha_0$-edge; hence no
$\alpha_0$-edge joins that pair to $Z_{\gamma(\delta)}$. Closedness of
$V(T)$ excludes an $\alpha_0$-edge from $V(T)$ to its exterior. Thus no
$\alpha_0$-edge leaves $Z_{\gamma(\delta)}$, while the only
$\delta$-edge from that set to the exterior of $V(T)$ is
$f^0_\delta$. It follows that the initial segment
$y_0P_\delta v^0_\delta$ is contained in $Z_{\gamma(\delta)}$ and, since
$v^0_\delta$ is the last vertex of $V(T)$ on $P_\delta$, the path
$P_\delta$ avoids $U_{\gamma(\delta)}$ altogether.

Consider the distinct $(\alpha_0,\delta)$-chain through
$u^1_{\gamma(\delta)}$. Equations \eqref{L41-3} and \eqref{L41-4} show
that it meets $V(T)$ only in $U_{\gamma(\delta)}$; if it later entered
$Z_{\gamma(\delta)}$, it would contain $f^0_\delta$ and hence coincide
with $P_\delta$. It contains no edge of the Tashkinov-tree sequence, since
the only tree edges incident with this pair have colour
$\gamma(\delta)$. Switching the chain therefore leaves $T$ Tashkinov: a
tree edge of colour $\alpha_0$, if any, retains its unchanged witness
$y_0$, and $\delta$ is unused on $T$. Since $\delta$ is defective and only
one $\delta$-boundary edge has its in-end in
$Z_{\gamma(\delta)}$, at least two such edges are incident with
$U_{\gamma(\delta)}$. There can be at most two, one at each vertex of
that pair, so there are exactly two. Together with the internal
$\alpha_0$-edge, they lie on the switched chain and become
$\alpha_0$-edges. Appending either one to $T$ is then possible because
$\alpha_0$ remains missing at $y_0$, contradicting the definition of
$t(G)$. This proves \eqref{L41-7}.

Besides $f^0_\delta$, denote the other two edges counted in
\eqref{L41-7} by
\[
 f^i_\delta=v^i_\delta u^i_\delta\quad(i=1,2),
 \qquad
 v^i_\delta\in Z_{\gamma(\delta)},
 \quad u^i_\delta\notin V(T),
\]
and put $U_\delta=\set{u^1_\delta,u^2_\delta}$. Neither
$u^1_\delta$ nor $u^2_\delta$ belongs to $P'_\delta$: otherwise its
unique incident $\delta$-edge $f^i_\delta$ would lie on that chain, forcing
the tree vertex $v^i_\delta\ne v^0_\delta$ onto it, contrary to
\eqref{L41-5}. Consequently,
\begin{equation}
 \varphi_\delta(f)=\varphi(f)
 \quad\text{for every edge $f$ incident with $U_\delta$.}
 \tag{8}\label{L41-8}
\end{equation}
The sequence
\[
 T_\delta=(T-\gamma(\delta),
            f^1_\delta,u^1_\delta,
            f^2_\delta,u^2_\delta)
\]
is a Tashkinov tree with respect to $e,\varphi_\delta$: both appended
edges have colour $\delta$, which is missing at $v^0_\delta$. It has
$t(G)$ vertices, so
\begin{equation}
 (T_\delta,e,\varphi_\delta)\in\Tcal(G).
 \tag{9}\label{L41-9}
\end{equation}

Closedness of this new maximum tree first gives
\begin{equation}
 E_G(u^1_\delta,u^2_\delta)\cap E_\alpha(e,\varphi)
 \ne\varnothing
 \quad\text{for every }
 \alpha\in\overline\varphi(Z_{\gamma(\delta)})
              \comp\set{\gamma(\delta)}.
 \tag{10}\label{L41-10}
\end{equation}
Indeed, $\alpha$ is not one of the two colours changed on $P'_\delta$:
$\delta$ is present on $Z_{\gamma(\delta)}$, and
$\gamma(\delta)$ is excluded. Thus $\alpha$ has the same edges and the
same missing vertex in $Z_{\gamma(\delta)}$ under both colourings. The
original closed set $V(T)$ has no $\alpha$-edge from
$Z_{\gamma(\delta)}$ to $U_\delta$. In the closed maximum tree
$T_\delta$, the $\alpha$-edge at either vertex of $U_\delta$ can neither
leave $V(T_\delta)$ nor end in $Z_{\gamma(\delta)}$. Since $\alpha$ is
missing elsewhere in $V(T_\delta)$, it is present at both vertices of
$U_\delta$, and their $\alpha$-edge must join them. Equation
\eqref{L41-8} transfers the conclusion back to $\varphi$.

We shall also need the two symmetric internal-edge conclusions
\begin{align}
 E_G(u^1_{\gamma(\delta)},u^2_{\gamma(\delta)})
   \cap E_\delta(e,\varphi)&\ne\varnothing,
   \tag{11}\label{L41-11}\\
 E_G(u^1_\delta,u^2_\delta)
   \cap E_{\gamma(\delta)}(e,\varphi)&\ne\varnothing.
   \tag{12}\label{L41-12}
\end{align}
We prove \eqref{L41-11} first. The set
$Z_{\gamma(\delta)}$ has $p+2s-2$ vertices and contains both ends of $e$.
By \Cref{maximum-triple-facts}(f) and elementarity,
\[
 |\overline\varphi(Z_{\gamma(\delta)})|\ge p+2s.
\]
At most $p-2+s$ colours are used on $T$: at most $p-2$ occur on its
initial coloured edges and at most one new colour occurs for each of the
$s$ blocks. Every such colour is missing on $Y\subseteq
Z_{\gamma(\delta)}$. Hence there is a colour
\[
 \beta\in
 \overline\varphi(Z_{\gamma(\delta)})
 \cap\Fcol(T,e,\varphi).
\]
Let $v$ be the vertex missing $\beta$ and consider
$P=P_v(\beta,\delta,\varphi)$. By \Cref{defective-chain}(b) and
\eqref{L41-7}, the chain contains $f^0_\delta,f^1_\delta,f^2_\delta$.
Equation \eqref{L41-10} gives an internal $\beta$-edge of $U_\delta$;
therefore $f^1_\delta$ and $f^2_\delta$ form one excursion of $P$ from
$Z_{\gamma(\delta)}$ and back to it.

Let $v'$ be the last vertex of $V(T)$ on $P$, in the order beginning at
$v$. It is an exit vertex, hence lies in $Y$ by \eqref{L41-2}. The final
edge leaving $V(T)$ must be one of
$f^0_\delta,f^1_\delta,f^2_\delta$. The latter two are paired by the
internal $\beta$-edge and form a returning excursion; neither can be the
final departure. Thus $v'=v^0_\delta$, and the final departure is
$f^0_\delta$.

The component $P$ cannot contain a vertex of
$U_{\gamma(\delta)}$. To see this, note first that \eqref{L41-3} gives an
internal $\beta$-edge in that pair, so no $\beta$-edge joins it to
$Z_{\gamma(\delta)}$; equation \eqref{L41-4} excludes a $\delta$-edge
between the two sets. Any route from $Z_{\gamma(\delta)}$ to the pair would
therefore have to leave $V(T)$ first. The departures through
$f^1_\delta,f^2_\delta$ are trapped in the already described excursion
through $U_\delta$, while the departure through $f^0_\delta$ occurs after
the last tree vertex. Hence no such route is possible. By
\Cref{defective-chain}(b), every $\delta$-coloured boundary edge of
$V(T)$ lies on $P$, so none is incident with
$U_{\gamma(\delta)}$. Together with \eqref{L41-4} and the fact that
$\delta$ is present at both vertices of that pair, this forces the internal
$\delta$-edge asserted in \eqref{L41-11}.

For \eqref{L41-12}, work with the globally maximum triple
$(T_\delta,e,\varphi_\delta)$. The same missing-colour count on
$Z_{\gamma(\delta)}$ gives a free colour $\beta$ for this triple. Under
$\varphi_\delta$, the colour $\gamma(\delta)$ is defective, and its
boundary edges with in-end in $Z_{\gamma(\delta)}$ are exactly
\[
 f^0_\delta,
 \quad f^1_{\gamma(\delta)},
 \quad f^2_{\gamma(\delta)}.
\]
The first acquired colour $\gamma(\delta)$ in \eqref{L41-6}; the other two
are the deleted block edges. There is no further one: before the switch,
closedness of $V(T)$ excluded every $\gamma(\delta)$-edge from
$Z_{\gamma(\delta)}$ to the exterior of $V(T)$, and the two block vertices
already use their unique $\gamma(\delta)$-edges.

Every vertex of $U_\delta$ misses a free colour for the new triple. Indeed,
each misses at least one colour, while every colour used on $T_\delta$ has
a missing witness in $Z_{\gamma(\delta)}$ and hence, by elementarity, is
present on $U_\delta$. Thus \Cref{exit-no-free} says that every exit vertex
of $T_\delta$ lies in $Z_{\gamma(\delta)}$. Now repeat the proof of
\eqref{L41-11}. The edges
$f^1_{\gamma(\delta)},f^2_{\gamma(\delta)}$ form a returning excursion
through $U_{\gamma(\delta)}$, because \eqref{L41-3} supplies its internal
$\beta$-edge. The pair $U_\delta$ has its own internal $\beta$-edge by
\eqref{L41-10}, and no $\gamma(\delta)$-edge joins it to
$Z_{\gamma(\delta)}$, by \eqref{L41-8} and the original closedness of
$V(T)$. Thus the same last-exit argument shows that the relevant
$(\beta,\gamma(\delta))$-chain avoids $U_\delta$. Consequently no
$\gamma(\delta)$-coloured boundary edge of $T_\delta$ is incident with
that pair. Since the colour is present at both of its vertices, they are
joined by a $\gamma(\delta)$-edge under $\varphi_\delta$, and hence under
$\varphi$ by \eqref{L41-8}. This proves \eqref{L41-12}.

Two further symmetric consequences will be used repeatedly:
\begin{align}
 E_G(u^1_{\gamma(\delta)},u^2_{\gamma(\delta)})
   \cap E_\alpha(e,\varphi)&\ne\varnothing
 &&\text{for every }\alpha\in\overline\varphi(U_\delta),
 \tag{13}\label{L41-13}\\
 E_G(u^1_\delta,u^2_\delta)
   \cap E_\alpha(e,\varphi)&\ne\varnothing
 &&\text{for every }\alpha\in
          \overline\varphi(U_{\gamma(\delta)}).
 \tag{14}\label{L41-14}
\end{align}
For \eqref{L41-13}, let $\alpha\in\overline\varphi(U_\delta)$. The
$\delta$-edges $f^1_\delta,f^2_\delta$ and \eqref{L41-12} show that
$\alpha\notin\set{\delta,\gamma(\delta)}$. Thus the
$\alpha$-edges are the same in $\varphi$ and $\varphi_\delta$.
Furthermore $\alpha\ne\alpha_0$, because $V(T_\delta)$ is elementary,
$\alpha$ is missing on $U_\delta$, and $\alpha_0$ is missing at $y_0$.
The $(\alpha_0,\alpha)$-chain through
$u^1_{\gamma(\delta)}$ is disjoint from $V(T_\delta)$: both colours are
missing on that closed set. It therefore meets $V(T)$ only inside
$U_{\gamma(\delta)}$ and contains no edge of the tree sequence $T$.
Switching it preserves $T$, since every colour used on a tree edge has an
unchanged missing witness in $Y$, which the chain avoids.

By \eqref{L41-3}, the chain contains both vertices of
$U_{\gamma(\delta)}$ through their internal $\alpha_0$-edge. If an
$\alpha$-edge joined the pair to the exterior of $V(T)$, the switch would
turn it into an $\alpha_0$-edge appendable to $T$. If an $\alpha$-edge
joined the pair to $Z_{\gamma(\delta)}$, it would be a boundary edge of the
closed maximum tree $T_\delta$ in a colour missing on that tree. Both are
impossible. Hence no $\alpha$-edge leaves the pair. Without an internal
$\alpha$-edge, both vertices of the pair would miss $\alpha$, contrary to
the elementarity of $V(T)$. This proves \eqref{L41-13}.

For \eqref{L41-14}, let
$\alpha\in\overline\varphi(U_{\gamma(\delta)})$. The block edges of colour
$\gamma(\delta)$ and \eqref{L41-11} give
$\alpha\notin\set{\gamma(\delta),\delta}$. The
$(\alpha_0,\alpha)$-chain through $u^1_\delta$ is disjoint from $V(T)$,
because both colours are missing on that closed set, and it meets
$V(T_\delta)$ only in $U_\delta$. Equation \eqref{L41-10} gives its
internal $\alpha_0$-edge in $U_\delta$. Switching the chain preserves
$T_\delta$: every initial tree-edge colour has a witness in $Y$, and the
two appended $\delta$-edges have the witness $v^0_\delta\in Y$; all these
vertices are avoided. The same switch-and-append argument, with
$T$ and $T_\delta$ interchanged, excludes every $\alpha$-edge leaving
$U_\delta$, and elementarity of $V(T_\delta)$ forces an internal one. This
proves \eqref{L41-14}.

We now prove the absorbing assertion
\begin{equation}
 U_\delta\subseteq A(T,e,\varphi)
 \quad\text{for every }\delta\in\Dcol(T,e,\varphi).
 \tag{15}\label{L41-15}
\end{equation}
Fix $u\in U_\delta$, let $\alpha\in\overline\varphi(u)$, and let
$\beta\in\Fcol(T,e,\varphi)\comp\set\alpha$. By \eqref{L41-12} and the
incident $\delta$-edge, $\alpha\notin\set{\delta,\gamma(\delta)}$; a free
colour $\beta$ is also distinct from those two colours. Hence the
$(\alpha,\beta)$-chain through $u$ is identical under $\varphi$ and
$\varphi_\delta$. Let $v\in V(T)$ be the unique vertex missing $\beta$.
If $v\in Z_{\gamma(\delta)}$, then $u,v\in V(T_\delta)$, and
\Cref{Scheide-Theorem-3.3}(e), applied to the maximum tree $T_\delta$,
shows that the chain is a path with ends $u,v$.

Suppose that $v\in U_{\gamma(\delta)}$. Write
$U_\delta=\set{u,u'}$ and
$U_{\gamma(\delta)}=\set{v,v'}$. Equation \eqref{L41-14} supplies the
internal $\beta$-edge $uu'$. Since $\alpha$ is missing at $u$ and
$V(T_\delta)$ is elementary and closed, the unique $\alpha$-edge at $u'$
joins $u'$ to a vertex of $Z_{\gamma(\delta)}$. Similarly, $\beta$ is
missing at $v$, and the elementarity and closedness of $V(T)$ imply that
the unique $\beta$-edge at $v'$ joins $v'$ to
$Z_{\gamma(\delta)}$. Finally, \eqref{L41-13} supplies the internal
$\alpha$-edge $v'v$.

These two edges entering $Z_{\gamma(\delta)}$ belong to the same
$(\alpha,\beta)$-component. Indeed, every vertex of
$Z_{\gamma(\delta)}$ is incident with both colours: $\alpha$ is missing at
$u\in V(T_\delta)$ and $\beta$ is missing at $v\in V(T)$. The closedness
of $V(T_\delta)$ excludes every $\alpha$-edge from
$Z_{\gamma(\delta)}$ to $U_{\gamma(\delta)}$ or to the exterior of
$V(T_\delta)$; because $u$ misses $\alpha$, the edge from $u'$ is the
unique $\alpha$-edge from $Z_{\gamma(\delta)}$ to $U_\delta$.
Similarly, the closedness of $V(T)$ excludes every $\beta$-edge from
$Z_{\gamma(\delta)}$ to $U_\delta$ or to the exterior of $V(T)$; because
$v$ misses $\beta$, the edge to $v'$ is the unique $\beta$-edge from
$Z_{\gamma(\delta)}$ to $U_{\gamma(\delta)}$. Thus these are exactly
the two boundary edges of $Z_{\gamma(\delta)}$ in the
$(\alpha,\beta)$-subgraph. Every vertex inside the set has degree two in
that subgraph, so its alternating paths pair the two boundary edges.
Consequently the chain beginning $u-\beta-u'-\alpha$ passes through
$Z_{\gamma(\delta)}$, then through $v'-\alpha-v$, and has ends $u,v$.
This is exactly the absorbing condition, so \eqref{L41-15} follows.

Set
\[
 X=V(T)\cup
   \bigcup_{\delta\in\Dcol(T,e,\varphi)}U_\delta.
\]
By \eqref{L41-15} and \Cref{absorbing-elementary},
\begin{equation}
 X\text{ is elementary with respect to }\varphi.
 \tag{16}\label{L41-16}
\end{equation}
For a fixed defective colour $\delta$, put
$X_\delta=V(T)\cup U_\delta$. We claim
\begin{equation}
 X_\delta\text{ is closed with respect to }\varphi.
 \tag{17}\label{L41-17}
\end{equation}
Let $\alpha\in\overline\varphi(X_\delta)$. Suppose first that $\alpha$ is
missing on $V(T)$. The closedness of $V(T)$ excludes an $\alpha$-edge from
that set to its exterior. The two vertices of $U_\delta$ are joined by an
$\alpha$-edge: use \eqref{L41-12} if
$\alpha=\gamma(\delta)$, use \eqref{L41-10} if its missing vertex lies in
$Z_{\gamma(\delta)}$ and $\alpha\ne\gamma(\delta)$, and use
\eqref{L41-14} if its missing vertex lies in
$U_{\gamma(\delta)}$. Thus no $\alpha$-edge leaves $X_\delta$.

Suppose instead that $\alpha$ is missing on $U_\delta$. The incident
$\delta$-edges and \eqref{L41-12} imply
$\alpha\notin\set{\delta,\gamma(\delta)}$. Hence the colour is unchanged
in $\varphi_\delta$, and it is missing on the closed maximum tree
$T_\delta$. No $\alpha$-edge can leave
$Z_{\gamma(\delta)}\cup U_\delta$. Equation \eqref{L41-13} gives an
internal $\alpha$-edge of the remaining pair
$U_{\gamma(\delta)}$, so no $\alpha$-edge leaves that pair either. This
proves \eqref{L41-17}.

Equation \eqref{L41-10}, applied to $\alpha_0$, supplies an internal
$\alpha_0$-edge in every $U_\delta$. Since a vertex is incident with at
most one edge of a given colour,
\begin{equation}
 U_\delta\text{ and }U_{\delta'}
 \text{ are either equal or disjoint.}
 \tag{18}\label{L41-18}
\end{equation}
We now prove
\begin{equation}
 X\text{ is closed with respect to }\varphi.
 \tag{19}\label{L41-19}
\end{equation}
Suppose that an edge $f$ of colour $\alpha$ leaves $X$, where
$\alpha\in\overline\varphi(X)$. Choose $\delta$ so that the in-end of $f$
lies in $X_\delta$. Since $X_\delta$ is closed, the vertex of $X$ missing
$\alpha$ is not in $X_\delta$; by \eqref{L41-18}, it lies in a pair
$U_{\delta'}$ disjoint from $U_\delta$. Closedness of $X_{\delta'}$ then
shows that the in-end of $f$ cannot lie in $V(T)$, so it lies in
$U_\delta$. Denote this end by $u$ and the outside end by $w$.

The $(\alpha_0,\alpha)$-chain through $u$ avoids $X_{\delta'}$, because
both colours are missing on that closed set. In particular it avoids
$V(T)$. By \eqref{L41-16}, the colour $\alpha$ is not in $\Gamma$:
every colour of $\Gamma$ is already missing on $V(T)$. Also
$\alpha\ne\delta$, since the $\delta$-edge $f^1_\delta$ crosses the
boundary of the closed set $X_{\delta'}$. We already have
$\alpha_0\notin\Gamma$ and $\alpha_0\ne\delta$. Therefore this chain is
the same under $\varphi$ and $\varphi_\delta$, and it contains no edge of
$T_\delta$. Switching it preserves all Tashkinov conditions for
$T_\delta$: every tree-edge colour has a missing witness in $V(T)$, which
the chain avoids. The edge $f$ lies on the chain and acquires colour
$\alpha_0$, while $\alpha_0$ remains missing at $y_0$. Appending $f,w$ to
$T_\delta$ contradicts the definition of $t(G)$. This proves
\eqref{L41-19}.

We require one final boundary fact. If $\alpha$ is present at every vertex
of $X$ and
\[
 P=P_{y_0}(\alpha_0,\alpha,\varphi),
\]
then
\begin{equation}
 |E(P)\cap E_G(X,V(G)\comp X)|=1.
 \tag{20}\label{L41-20}
\end{equation}
The path has one end $y_0$ and its other end outside $X$, because $X$ is
elementary, $\alpha_0$ is missing at $y_0$, and $\alpha$ is missing at no
vertex of $X$. Let $v$ be the last vertex of $X$ on the path, and let $g$
be the following edge. By \eqref{L41-19}, no $\alpha_0$-edge leaves $X$,
so $\varphi(g)=\alpha$. We show that $y_0Pv\subseteq X$.

If $v\in V(T)$ and $\alpha$ is not defective for $T$, then $g$ is the
unique $\alpha$-edge leaving $V(T)$, while no $\alpha_0$-edge leaves that
set. Thus $y_0Pv\subseteq V(T)$. If $v\in V(T)$ and $\alpha$ is defective,
then \eqref{L41-7} and \eqref{L41-11}, with $\delta=\alpha$, show that the
only $\alpha$-edge leaving $X_\alpha$ is $f^0_\alpha$: the other two
boundary edges of $V(T)$ end in $U_\alpha$, and no such edge is incident
with $U_{\gamma(\alpha)}$. Since $X_\alpha$ is closed and
$\alpha_0$ is missing there, the initial segment lies in $X_\alpha$.

Finally suppose that $v\in U_\delta$. The edge $g$ is different from the
incident $\delta$-edge and from the internal
$\gamma(\delta)$-edge in \eqref{L41-12}; hence
$\alpha\notin\set{\delta,\gamma(\delta)}$. The path is therefore unchanged
under $\varphi_\delta$, and $v$ is its last vertex in
$V(T_\delta)\subseteq X$. Every colour used on $T_\delta$ has a missing witness in
$Z_{\gamma(\delta)}$: this is clear for the inherited tree edges, and the
two appended $\delta$-edges use $v^0_\delta$ as their witness. Since
$v\in U_\delta$, elementarity therefore prevents any used colour from
being missing at $v$. Thus every colour missing at $v$ is unused, and $v$
misses a free colour. If $\alpha$ were defective for
$T_\delta$, the vertex $v$ would be an exit vertex for the chain from
$y_0$ in the colours $\alpha_0,\alpha$, contrary to
\Cref{exit-no-free}. Hence $\alpha$ is not defective for $T_\delta$.
The edge $g$ is its unique $\alpha$-boundary edge, and no
$\alpha_0$-edge leaves the closed tree. Therefore
$y_0Pv\subseteq V(T_\delta)$. All cases prove \eqref{L41-20}.

We finish by proving
\begin{equation}
 X\text{ is strongly closed with respect to }\varphi.
 \tag{21}\label{L41-21}
\end{equation}
Suppose not. By \eqref{L41-19}, some colour $\alpha$ that is present at
every vertex of $X$ occurs on at least two edges of $\partial(X)$. The set
$X$ has odd order: $V(T)$ has odd order and, by \eqref{L41-18}, the union
adds disjoint two-vertex pairs. Counting incidences of the
$\alpha$-matching therefore shows that an odd number, hence at least three,
of its edges leave $X$.

At most one of these boundary edges has its in-end in $V(T)$. Otherwise
$\alpha$ would be defective for $T$, while \eqref{L41-7} and
\eqref{L41-11} would show that only $f^0_\alpha$ among its boundary edges
can leave $X_\alpha$. Thus at least two boundary edges have their in-end in
$X\comp V(T)$.

By \eqref{L41-20}, the path
$P_{y_0}(\alpha_0,\alpha,\varphi)$ uses exactly one boundary edge of $X$.
Choose another one, $f=wu$, where $w\in U_\delta$ and $u\notin X$, not on
that path. The chain
$Q=P_u(\alpha_0,\alpha,\varphi)$ is disjoint from the first chain. It also
avoids $V(T)$. If $\alpha$ is defective for $T$, the closed set
$X_\alpha$ has only the edge $f^0_\alpha$ of colour $\alpha$ leaving it,
and that edge lies on the first chain; hence the disjoint chain $Q$ cannot
enter $X_\alpha\supseteq V(T)$. If $\alpha$ is not defective, the unique
$\alpha$-boundary edge of $V(T)$ lies on the first chain, and again $Q$
cannot enter $V(T)$. No $\alpha_0$-edge leaves $V(T)$ in either case.

The edge $f$ is distinct from the incident $\delta$-edge at $w$ and from
the internal $\gamma(\delta)$-edge supplied by \eqref{L41-12}. Hence
$\alpha\notin\set{\delta,\gamma(\delta)}$; the same is true of
$\alpha_0$. Thus $Q$ is unchanged under $\varphi_\delta$ and contains no
edge of $T_\delta$. Switching it preserves $T_\delta$, because all
missing witnesses for its tree-edge colours lie in the avoided set
$V(T)$. The edge $f$ then has colour $\alpha_0$, while $\alpha_0$ remains
missing at $y_0$. Appending $f,u$ exceeds $t(G)$, a contradiction. This
proves \eqref{L41-21}.

By \eqref{L41-16} and \eqref{L41-21}, the set $X$ is elementary and
strongly closed, and it contains both ends of $e$. The implication \textup{(iv)} $\Rightarrow$ \textup{(i)} in
\Cref{egraph2} now proves that $G$ is elementary.
\end{proof}

\begin{proposition}[Source Proposition~4.2]
\label{Scheide-Proposition-4.2}\label{large-height}
Let $G$ be critical with $\chi'(G)=k+1$ and $k\ge\Delta(G)+1$. If
$h(G)>t(G)-4$, then $G$ is elementary.
\end{proposition}

\begin{proof}
Both $h(G)$ and $t(G)$ are odd, by \Cref{balanced-triples}(a) and
\Cref{maximum-triple-facts}(a). Hence the hypothesis gives
$h(G)=t(G)$ or $h(G)=t(G)-2$. Choose a balanced triple
$(T,e,\varphi)$.

Suppose first that $h(G)=t(G)$. Then $T$ consists of its trunk, and only the
two trunk colours $\alpha,\beta$ are used on it. Every vertex of $V(T)$
misses a free colour. Indeed, the two ends of $e$ each miss at least two
colours and at most one of $\alpha,\beta$ at that vertex, while every other
vertex misses at least one colour and, by elementarity, cannot miss either
trunk colour, which is already missing at one of the first two vertices. If
a defective colour existed, \Cref{defective-chain}(d) would give an exit
vertex missing no free colour, a contradiction. Thus there is no defective
colour. By \Cref{maximum-triple-facts}(b),(c), $V(T)$ is elementary and
strongly closed, and \Cref{egraph2} proves that $G$ is elementary.

Suppose now that $h(G)=t(G)-2$. Write $p=h(G)$. The balanced tree has the
form
\[
 T=(y_0,e_1,y_1,\ldots,e_{p-1},y_{p-1},f_1,u_1,f_2,u_2),
\]
where its trunk uses the two colours $\alpha,\beta$ and
\[
 \varphi(f_1)=\varphi(f_2)=\gamma
 \in\overline\varphi(\set{y_0,\ldots,y_{p-1}}).
\]
The colour $\gamma$ is distinct from $\alpha,\beta$: the closing
$(\alpha,\beta)$-path is the full two-coloured component of the ends of
$e$, so no $\alpha$- or $\beta$-edge can leave its vertex set along the
tail. Properness also shows that both $f_1$ and $f_2$ join their new
vertices to the trunk; the second cannot use $u_1$, already incident with
the $\gamma$-edge $f_1$.

Let $y_r$ be the unique trunk vertex missing $\gamma$. Every vertex of
$V(T)\comp\set{y_r}$ misses a free colour: all used colours are
$\alpha,\beta,\gamma$, and their missing vertices are already prescribed;
the ends of $e$ have an additional missing colour when needed. By
\Cref{exit-no-free}, $F(T,e,\varphi)\subseteq\set{y_r}$. The hypotheses of
\Cref{structural-lemma} hold with $s=1$, $Y$ equal to the trunk, and
$\Gamma=\set{\gamma}$. That lemma proves that $G$ is elementary.
\end{proof}

\begin{proposition}[Source Proposition~4.3]
\label{Scheide-Proposition-4.3}\label{small-height}
Let $G$ be critical with $\chi'(G)=k+1$ and $k\ge\Delta(G)+1$. If
$h(G)<5$, then $G$ is elementary.
\end{proposition}

\begin{proof}
By \Cref{balanced-triples}(a), $h(G)=3$. Choose a balanced triple and,
after naming its trunk colours, write
\[
 T=(y_0,e_1,y_1,e_2,y_2,e_3,y_3,e_4,y_4,\ldots),
\]
where $(y_0,e_1,y_1,e_2,y_2)$ is the trunk,
$\alpha\in\overline\varphi(y_0)$,
$\beta\in\overline\varphi(y_1)$, and a closing edge
$q\in E_G(y_2,y_0)$ has colour $\beta$.

Assume that $G$ is not elementary. By \Cref{large-height},
$t(G)\ge h(G)+4=7$, so the displayed pair $e_3,e_4$ exists. Its common
colour $\gamma$ is missing at a trunk vertex, is distinct from the trunk
colours, and both edges join $\set{y_1,y_2}$ to
$\set{y_3,y_4}$. Rotate the balanced trunk so that the vertex missing
$\gamma$ becomes $y_0$, and rename the two trunk colours if necessary.
Then $y_0,y_1$ form a $(\gamma,\beta)$-pair.

By \Cref{Scheide-Theorem-3.3}(e), their
$(\gamma,\beta)$-chain is an even path inside $V(T)$. Starting at $y_0$,
it uses the closing $\beta$-edge $q$ to $y_2$ and then the unique
$\gamma$-edge at $y_2$. At the other end it uses the unique
$\gamma$-edge at $y_1$. The two such edges are exactly $e_3,e_4$, one at
$y_1$ and one at $y_2$, so the path contains all five vertices
$y_0,\ldots,y_4$. By \Cref{normalisation}, there is a normal globally
maximum Tashkinov tree of height at least five, contradicting $h(G)=3$.
\end{proof}

\begin{lemma}[Height-five configuration; source Lemma~4.4]
\label{Scheide-Lemma-4.4}\label{height-five}
Let $G$ be critical with $\chi'(G)=k+1$, $k\ge\Delta(G)+1$, and
$h(G)=5$. Let $(T,e,\varphi)\in\Tcal^N(G)$ have trunk
\[
 T_0=(y_0,e_1,y_1,e_2,y_2,e_3,y_3,e_4,y_4),
\]
with colours $\alpha\in\overline\varphi(y_0)$ and
$\beta\in\overline\varphi(y_1)$. Suppose that
$\gamma\in\overline\varphi(y_0)$ and
\[
 E_G(V(T_0),V(T)\comp V(T_0))\cap E_\gamma(e,\varphi)
 \ne\varnothing.
\]
Then:
\begin{enumerate}[label=\textup{(\arabic*)}]
\item there are edges
\[
 f_1\in E_G(y_1,V(T)\comp V(T_0)),\quad
 f_2\in E_G(y_4,V(T)\comp V(T_0)),\quad
 f_3\in E_G(y_2,y_3)
\]
all coloured $\gamma$;
\item if $v_1,v_2$ are the ends of $f_1,f_2$ outside $V(T_0)$, then
      $E_G(v_1,v_2)$ contains both an $\alpha$-edge and a $\beta$-edge.
\end{enumerate}
\end{lemma}

\begin{proof}
By the definition of the trunk,
\[
 e_1=e,
 \qquad
 \varphi(e_2)=\varphi(e_4)=\alpha,
 \qquad
 \varphi(e_3)=\beta,
\]
and there is a closing edge $q\in E_G(y_4,y_0)$ with
$\varphi(q)=\beta$.

Let $P_1=P_{y_0}(\gamma,\beta,\varphi)$. By
\Cref{Scheide-Theorem-3.3}(e), it is an even path with ends $y_0,y_1$ and
vertices in $V(T)$. Applying \Cref{normalisation} to this pair shows that
$|V(P_1)|\le h(G)=5$. The edge $q$ lies on $P_1$, so
$y_0,y_1,y_4\in V(P_1)$.

No $\gamma$-edge joins $\set{y_2,y_3}$ to
$V(T)\comp V(T_0)$. Suppose that such an edge $g$ exists and let $v$ be
its end outside $V(T_0)$. If one of $y_2,y_3,v$ lay on $P_1$, then the
$\gamma$-edge $g$ and the $\beta$-edge $e_3$ would put all three on
$P_1$, producing at least six vertices together with $y_0,y_1,y_4$.
Thus all three avoid $P_1$.

Switch $P_1$. The full tree $T$ remains Tashkinov by
\Cref{Scheide-Theorem-3.3}(e), because $\gamma$ is unused on the initial
two-vertex tree. In the new colouring $\gamma$ is missing at $y_1$, the
closing edge $q$ has colour $\gamma$, and $g$ is unchanged. The
$(\alpha,\gamma)$-chain between $y_0$ and $y_1$ contains the forced
segments
\[
 y_0,q,y_4,e_4,y_3
 \qquad\text{and}\qquad
 y_1,e_2,y_2,
\]
and it also contains $g$ and $v$. Hence it has at least six vertices.
Normalisation would give a globally maximum normal tree of height greater
than five, a contradiction.

The $\gamma$-edges at $y_2,y_3$ therefore stay in $V(T_0)$. If they were
distinct, matching and the fact that $\gamma$ is missing at $y_0$ would
force them to end at $y_1,y_4$, leaving no $\gamma$-edge from the trunk to
its complement. Hence they coincide in an edge
$f_3\in E_G(y_2,y_3)$. The $\gamma$-edges at $y_1,y_4$ cannot coincide,
for the same reason, and their other ends cannot be elsewhere in the trunk.
They are the edges $f_1,f_2$ asserted in (1).

Let their outside ends be $v_1,v_2$. The path $P_1$ contains the two
segments $y_1,f_1,v_1$ and $y_0,q,y_4,f_2,v_2$. These comprise five
distinct vertices, and $|V(P_1)|\le5$. Therefore a $\beta$-edge joins
$v_1$ and $v_2$ and completes the path.

Finally switch the original $(\alpha,\beta)$-trunk path
\[
 P(y_1,e_2,y_2,e_3,y_3,e_4,y_4,q,y_0).
\]
Again $T$ remains Tashkinov by \Cref{Scheide-Theorem-3.3}(e). Now
$\alpha$ is missing at $y_1$, $\gamma$ is still missing at $y_0$, and
$q$ has colour $\alpha$. The $(\gamma,\alpha)$-chain between $y_0,y_1$
contains the five vertices in the two segments
$y_0,q,y_4,f_2,v_2$ and $y_1,f_1,v_1$. By maximality of $h(G)$ it has
exactly five vertices, so an $\alpha$-edge joins $v_1,v_2$. This proves
(2).
\end{proof}

\begin{proposition}[Small Tashkinov order; source Proposition~4.5;
 target label \target{ThmScheide}]
\label{Scheide-Proposition-4.5}\label{ThmScheide}
Let $G$ be a critical multigraph with $\chi'(G)\ge\Delta(G)+2$. If
$t(G)<11$, then $G$ is elementary.
\end{proposition}

\begin{proof}
Put $k=\chi'(G)-1$, so $k\ge\Delta(G)+1$ and the preceding results apply.
Suppose for a contradiction that $G$ is not elementary. Since $t(G)$ is
odd, \Cref{large-height,small-height} force
\[
 t(G)=9,
 \qquad
 h(G)=5.
\]
Choose a balanced triple and write
\[
 T=(y_0,e_1,y_1,e_2,y_2,e_3,y_3,e_4,y_4,
     e_5,y_5,e_6,y_6,e_7,y_7,e_8,y_8),
\]
where $e_1=e$, the trunk colours are $\alpha,\beta$, and
\[
 \varphi(e_5)=\varphi(e_6)=\gamma_1,
 \qquad
 \varphi(e_7)=\varphi(e_8)=\gamma_2.
\]
The first tail colour $\gamma_1$ is missing at a trunk vertex. Rotate the
balanced trunk so that it is missing at $y_0$.

Apply \Cref{height-five} with $\gamma=\gamma_1$. There are vertices
$v_1,v_2\in\set{y_5,y_6,y_7,y_8}$ such that
\[
 y_1v_1,\qquad y_4v_2,
 \qquad y_2y_3
\]
have colour $\gamma_1$, and $v_1v_2$ has both an $\alpha$-edge and a
$\beta$-edge. The two tree edges $e_5,e_6$ are precisely the first two of
these $\gamma_1$-edges: their in-ends cannot be $y_0$, where
$\gamma_1$ is missing, or $y_2,y_3$, already incident with the internal
$\gamma_1$-edge. Hence
\[
 \set{v_1,v_2}=\set{y_5,y_6}.
\]
The two equal-coloured edges may be interchanged in their block, so we
reorder them, if necessary, and henceforth take $v_1=y_5$ and $v_2=y_6$.
In particular, every vertex among $y_0,\ldots,y_6$, except the unique
missing vertex for each of $\alpha,\beta,\gamma_1$, is already incident
with an edge of that colour. It follows at once that
\[
 \gamma_2\notin\set{\alpha,\beta,\gamma_1}:
\]
if equality held, neither of the two new vertices $y_7,y_8$ could be
joined by its tree edge to an earlier vertex not already incident with that
colour (and the unique missing vertex cannot be incident with such an
edge).

Since $V(T)$ is elementary and closed, the remaining unmatched incidences
of each of the colours $\gamma_1,\alpha,\beta$ force three edges between
$y_7$ and $y_8$, one of each colour. We may also assume, after replacing
and relabelling the balanced tree if necessary, that $\gamma_2$ is missing
on the trunk. Indeed, its missing witness precedes $e_7$, so it lies either
on the trunk or in $\set{y_5,y_6}$. In the latter case use the balanced
tree
\[
 (y_0,e_1,y_1,f_1,v_1,g_\beta,v_2,f_2,y_4,
   e_2,y_2,e_4,y_3,e_7,y_7,e_8,y_8),
\]
where $f_1=y_1v_1$ and $f_2=v_2y_4$ have colour $\gamma_1$, and
$g_\beta$ is the $\beta$-edge between $v_1,v_2$. Its trunk colours are
$\gamma_1,\beta$, its first tail pair has colour $\alpha$, and the vertex
missing $\gamma_2$ is now on the trunk. Relabel the vertices by their
positions in this new sequence as $y_0,\ldots,y_8$; then relabel the old
$\gamma_1$ as the new trunk colour $\alpha$ and the old $\alpha$ as the
new first-tail colour $\gamma_1$. The preceding conclusions, including
the three differently coloured edges between $y_7,y_8$, retain exactly the
same form.

We claim that a $\gamma_2$-edge joins the trunk to
$\set{y_7,y_8}$. Suppose not. Then the two tree edges $e_7,e_8$ join
$\set{y_5,y_6}$ to $\set{y_7,y_8}$, one at each vertex of each pair;
relabel so that $e_7=y_5y_7$ and $e_8=y_6y_8$. Together with the
$\beta$-edges $y_5y_6$ and $y_7y_8$, these form an alternating
$(\gamma_2,\beta)$-cycle. Switch that cycle. The displayed tree sequence
remains Tashkinov: $e_7,e_8$ now have colour $\beta$, which is missing at
a trunk vertex, and all other tree edges are unchanged. In the new
colouring the $(\gamma_1,\beta)$-path between the first two trunk vertices
contains, in order, the seven vertices
\[
 y_1,
 \quad v_1,
 \quad y_7,
 \quad y_8,
 \quad v_2,
 \quad y_4,
 \quad y_0,
\]
where $\set{v_1,v_2}=\set{y_5,y_6}$. By \Cref{normalisation}, this gives
height at least seven, contradicting $h(G)=5$. The claim follows.

Let $y_j$ be the trunk vertex missing $\gamma_2$. Temporarily rotate
the trunk so that $y_j$ is first and apply \Cref{height-five}. The rotation
changes only the $\alpha,\beta$ colours on the trunk cycle; it leaves every
$\gamma_2$-edge and every edge outside that cycle unchanged. The lemma gives
exactly two $\gamma_2$-edges from the trunk to its complement: the other two
trunk vertices incident with $\gamma_2$ are joined internally. The outside
ends of the two crossing edges are joined by an edge of each trunk colour,
in particular by an $\alpha$-edge. At least one outside end is $y_7$ or
$y_8$ by the preceding claim. But an $\alpha$-edge already joins $y_7$ and
$y_8$; properness therefore forces the two outside ends to be exactly
$y_7,y_8$. Returning to the original orientation and notation, the same two
$\gamma_2$-edges remain, so $e_7,e_8$ both join the trunk to these two
vertices.

We have now obtained precisely the configuration in
\Cref{structural-lemma}, with
\[
 Y=\set{y_0,\ldots,y_4},
 \qquad
 \Gamma=\set{\gamma_1,\gamma_2},
 \qquad s=2.
\]
The only colours used on $T$ are
$\alpha,\beta,\gamma_1,\gamma_2$. The trunk colours $\alpha,\beta$ are
missing at the first and second trunk vertices, respectively;
$\gamma_1$ is missing at $y_0$, and $\gamma_2$ is missing at $y_j$.
Every non-end of $e$ outside $\set{y_0,y_j}$ misses at least one colour,
and none of these four used colours can be missing there, by elementarity.
Each end of $e$ misses at least two colours; outside the same exceptional
set, at most one of the four used colours is missing there. Thus every
vertex outside $\set{y_0,y_j}$ misses an unused colour, hence a free
colour. Therefore \Cref{exit-no-free} gives
\[
 F(T,e,\varphi)\subseteq\set{y_0,y_j}.
\]
Each possible exit vertex in this set misses a colour of $\Gamma$, so (S4)
holds; (S1)--(S3) were established above. The structural lemma says that
$G$ is elementary, the final contradiction.
\end{proof}

\section{Dependency boundary and concordance}

The proof of \Cref{ThmScheide} uses no result from the portion of
Scheide's article beginning with Lemma~4.6.  Its source dependency cone is
exactly
\[
\begin{gathered}
 \text{Theorems 2.8, 3.2, 3.3; Propositions 3.4--3.6;}\\
 \text{Lemmas 3.8, 3.11, 3.12, 4.1, 4.4; Propositions 4.2, 4.3, 4.5}.
\end{gathered}
\]
Theorems~2.8 and~3.2 are the two results proved in the separately supplied
book extract, under the labels \texttt{thm:egraph2} and
\texttt{thm:tashkinov}. In the present note, Theorem~3.3,
Propositions~3.4--3.6, Lemmas~3.8, 3.11, 3.12, 4.1, 4.4, and
Propositions~4.2, 4.3, 4.5 are all proved. Thus these two files together
cover the entire displayed dependency cone. The later truncation argument
of source Lemma~4.7 and every result depending on it play no role here.

For source search and formalization, the endpoint has both labels
\target{Scheide-Proposition-4.5} and \target{ThmScheide}; the elementary-set
criterion has \target{Scheide-Theorem-2.8} and \target{egraph2}; and the
Tashkinov-tree theorem has \target{Scheide-Theorem-3.2} and
\target{TashTree}. The companion file additionally labels those proofs
\texttt{thm:egraph2} and \texttt{thm:tashkinov}. The definition of $t(G)$
carries the exact target label \target{torder}.

\begin{thebibliography}{99}

\bibitem[CJZ]{CJZ}
Guantao Chen, Guangming Jing, and Wenan Zang,
``Proof of the Goldberg--Seymour conjecture on edge-colorings of
multigraphs,''
arXiv:1901.10316v2, 7 June 2022.

\bibitem[FST]{FST}
Lene M. Favrholdt, Michael Stiebitz, and Bjarne Toft,
\emph{Graph Edge Colouring: Vizing's Theorem and Goldberg's Conjecture},
DMF-2006-10-003, IMADA-PP-2006-20, University of Southern Denmark,
preprint, 2006.

\bibitem[S]{S}
Diego Scheide,
``Graph edge coloring: Tashkinov trees and Goldberg's conjecture,''
\emph{Journal of Combinatorial Theory, Series B} \textbf{100} (2010),
68--96,
\href{https://doi.org/10.1016/j.jctb.2009.04.001}
     {doi:10.1016/j.jctb.2009.04.001}.

\bibitem[SSTF]{SSTF}
Michael Stiebitz, Diego Scheide, Bjarne Toft, and Lene M. Favrholdt,
\emph{Graph Edge Colouring: Vizing's Theorem and Goldberg's Conjecture},
Wiley, 2012.

\bibitem[T]{T}
V. A. Tashkinov,
``On an algorithm to colour the edges of a multigraph'' (in Russian),
\emph{Diskretnyi Analiz i Issledovanie Operatsii} \textbf{7} (2000),
72--85.

\end{thebibliography}

\end{document}
